\documentclass{problemset}
\usepackage{amsmath}

\usepackage{lipsum}
%\usepackage{showframe}
\usepackage{layout}


\usepackage[charter,cal=cmcal]{mathdesign} %different font
\usepackage{microtype}
\usepackage{amsmath}
%\usepackage{amsfonts}
%\usepackage{amssymb}
\usepackage{graphicx}
\usepackage[inline]{enumitem}
\usepackage{xparse}
\usepackage{ifthen}
\usepackage{graphicx}
\usepackage{caption}
\usepackage{subcaption}
\usepackage{kbordermatrix}
\usepackage{color}
\usepackage{tikz}
\usepackage{fancyhdr}
\usepackage{calc}
\usepackage[hidelinks]{hyperref}

\usepackage{pgfplots}
\pgfplotsset{compat=newest}
\usetikzlibrary{automata,arrows,positioning,calc}
%%%
% Useful Linear Algebra macros
%%%
\newcommand{\ul}{$\underline{\phantom{xxx}}$}
\newcommand{\ull}{\underline{\phantom{xxx}}}
\newcommand{\xh}{{\hat {\mathbf x}}}
\newcommand{\yh}{{\hat {\mathbf y}}}
\newcommand{\zh}{{\hat {\mathbf z}}}
\newcommand{\R}{\mathbb{R}}
\newcommand{\Z}{\mathbb{Z}}
\newcommand{\N}{\mathbb{N}}
\newcommand{\E}{\mathbb{E}}
\newcommand{\proj}{\mathrm{proj}}
\newcommand{\Proj}{\mathrm{proj}}
\newcommand{\Perp}{\mathrm{perp}}
\renewcommand{\span}{\mathrm{span}\,}
\newcommand{\Span}{\mathrm{span}\,}
\newcommand{\Img}{\mathrm{img}\,}
\newcommand{\Null}{\mathrm{null}\,}
\newcommand{\Range}{\mathrm{range}\,}
\newcommand{\rref}{\mathrm{rref}}
\newcommand{\rank}{\mathrm{rank}}
\newcommand{\Rank}{\mathrm{rank}}
\newcommand{\nnul}{\mathrm{nullity}}
\newcommand{\mat}[1]{\begin{bmatrix}#1\end{bmatrix}}
\newcommand{\chr}{\mathrm{char}}
\renewcommand{\d}{\mathrm{d}}
\renewcommand{\P}{\mathbb{P}}

\newcommand{\fatrule}[1]{\vspace{.3cm}\hrule {\hfill \sf #1}}

%\tcbuselibrary{skins}
%\usetikzlibrary{shadings}

%\definecolor{defcolor}{rgb}{.05,.4,.15}
%\colorlet{defframecolor}{green!50!black}
\declareoutlinedbox{lemma}{LEMMA}{LEM}{defframecolor}{defcolor}
\newenvironment{proof}{\emph{Proof.}}{\hfill$\square$}
%%%
% Set up the margins to use a fairly large area of the page
%%%
\textwidth=6in
\topmargin=-1in
\textheight=10in
\parskip=.07in
\parindent=0in

\begin{document}
%\pagestyle{empty}
\pagestyle{fancy}
\rfoot{\footnotesize\it \copyright\,Jason Siefken, 2016 \ \makebox(30,5){\includegraphics[height=1.2em]{by-sa.pdf}}}
\renewcommand{\headrulewidth}{0pt}

\begin{center}
	{\huge\bf Differential Equations \\{\sc Math 281-2}}\\

\vspace{.7in}
{
\it \copyright\,Jason Siefken, 2016 \\
Creative Commons By-Attribution Share-Alike\, \makebox(30,5){\includegraphics[height=1.2em]{by-sa.pdf}}
}
\end{center}

\section*{What is a Differential Equation?}


	A function $f:\R^n\to \R$ is a relationship between numbers.
	For example, if $f(x,y)=x^2-y^2$, the equation $f(x,y)=k$ has solutions which are pairs of numbers 
	$(x_0,y_0)$ satisfying $x_0^2-y_0^2=k$.

	Relationships between numbers are great, but we're graduating
	to relationships between functions.  Consider
	\[
		\frac{\d}{\d x}:\{\text{functions}\}\to\{\text{functions}\}.
	\]
	We might then consider $D(f) = \frac{\d f}{\d x}$ and ask for the set of solutions
	to
	\[
		D(f) = 0.
	\]
	From calculus you know that the only functions with zero derivative everywhere
	are constant functions.  Thus, $D(f)=0$ has solutions $f(x)=k$ for constant $k$.

	What about $D(f)=f$?  This equation has solutions $f(x)=0$ and $f(x)=ke^x$ for any
	constant $k$.

	Equation where the solutions are functions are called \emph{functional equations}.

	\begin{definition}[Differential Equation]
		A functional equation that involves derivatives is called
		a \emph{differential equation}.  An differential equation
		is called \emph{ordinary} if it is a differential equation
		without partial derivatives.
	\end{definition}

	We sometimes abbreviate \emph{ordinary differential equation} as ODE.

	In this class, we'll only deal with ordinary differential equations.

	Why is this idea useful?  Derivatives capture rates of change,
	and rates of change seem to relate to all sorts of quantities in nature. So, in a word,
	\emph{modeling}.

	\subsection*{Examples of Differential Equations}
	\begin{enumerate}
		\item From physics, if you're in a uniform gravitational field, the change
			in your velocity is constant with respect to time.  That is,
			\[
				\frac{\d v}{\d t} = a.
			\]
			You've discovered acceleration!
		\item Newton's law of cooling states that if $T(t)$ is the temperature
			of a body at time $T$ and $W$ is the ambient temperature, then
			\[
				\frac{ \d T}{\d t} = \alpha(T-W)
			\]
			for some $\alpha$.
		\item Imagine you're hungry.  Let $A(t)$ be your appetite at time $t$, and let
			$F(t)$ be how much food you've eaten.  You might imagine
			\[
				F'(t) = \alpha A(t). \qquad\qquad\text{(What would this mean?)}
			\]
			Now, if $T$ is the capacity of your tummy, you might also imagine
			\[
				A'(t) =\frac{-\beta}{T-F(t)}.\qquad\qquad\text{(Is this reasonable?)}
			\]
			This is a \emph{system} of differential equations.

			You might also suspect that it takes time for your appetite to change
			after you've eaten food.  We could model this with the equation
			\[
				A'(t) =\frac{-\beta}{T-F(t-k)}
			\]
			for some time delay $k$.
	\end{enumerate}

	You can dream up differential equations to model just about any situation.
	However, there is a cruel fact of life: \emph{almost all differential equations
	do not have elementary closed-form solutions}.  That is, almost
	every differential equation you dream up, you won't be able to write down
	a formula for its solution.

	Because of this, we will largely focus on modeling and estimating behavior of solutions
	and what we can do with solutions, rather than how to solve particular differential
	equations.

	\section*{Terminology}

	A differential equation (or system) is called $n$th order if $\frac{\d^n}{\d x^n}$ is the highest
	order derivative that appears.

	A differential equation is called \emph{linear} if it can be written in the form
	\[
		h_n(t)\frac{\d^n f(t)}{\d t^n} + h_{n-1}(t)\frac{\d^{n-1} f(t)}{\d t^{n-1}}
		+\cdots +h_1(t)\frac{\d f(t)}{\d t} + h_0(t) f(t) + h(t) = 0
	\]
	for some function $h_n,\ldots, h_1,h_0,h$ that depend only on $t$.  (It will turn out
	that linear differential equations with constant coefficients is one of the few types of
	differential equations we know how to solve.)  If $h(t)=0$ for some linear
	differential equation, it is called \emph{homogeneous}.

	\subsection*{Visualizing first order ODE's}

	Suppose we rewrite a first order ODE as
	\[
		f'(t) = F(t,f(t)).
	\]
	We can visualize this differential equation with a \emph{slope field}.  Letting $y=f(t)$,
	we can plot in the $ty$-plane little line segments representing the slope at each point.

	Let's visualize $f'(t) = f(t) + t$.

	\begin{center}
	{
	\def\length{sqrt(1+(x+y)^2)}
	\begin{tikzpicture}
	\begin{axis}[domain=-3:3, view={0}{90}, title={$\dfrac{\mathrm{d}f(t)}{\mathrm{d}t}=f(t)+t$}]
		\addplot3[blue, quiver={u={1/(\length)}, v={(x+y)/(\length)}, scale arrows=0.45}, -stealth,samples=12] {0};
		\draw (-3,0) -- (3,0);
		\draw (0,-3) -- (0,3);
	\end{axis}
	\end{tikzpicture}
	}
	\end{center}

	Looking at the slope field, we could imagine curves whose slope follows the pattern of the field
	a curve that \emph{flows} along the slope field).

	If we specify \emph{initial conditions}, like $f(-1)=0$, we've specified
	that $(-1,0)$ lies on our solution curve in the $ty$-plane.  From the picture, it looks like
	$f(t) = -t-1$ is the corresponding solutions, and checking, indeed this is a solution for
	the given initial conditions!

	Specifying different initial conditions, like $f(0)=0$, we expect to see a totally different solution curve.
	The solution for these initial conditions is $f(t)=e^{t-1}-t-1$.

	We can see from the slope field that depending on the initial condition, there are several different
	``looks'' a solution may have.  The line $y=-t-1$ seems to divide the set of initial
	conditions into solutions that curve up and those that curve down.

	Can we see this from the equation?

	\begin{quote}
		Consider the two cases, one where our initial conditions start above $f(t)=-t-1$
		and the other where they start below.

		$f < -t -1$.  In this case, $f+t<1$.  From the differential equation $f'=f+t$, we can see that that the slope
		starts out negative.  In particular $f'=\alpha <-1$.  If we increase $t$ by a small amount, $f$ will decrease 
		by approximately $\alpha \Delta t$.  Thus, our new slope will be $f_{\text{new}}'
		\approx f+t+\alpha \Delta t+\Delta t = f+t+(1+\alpha)\Delta t< -1$ since $1+\alpha < 0$.  Not only that, 
		$f_{\text{new}}' < f'$.  So, given these initial conditions, we'd expect the solution to head towards
		negative $y$ values at a faster and faster rate.

		$f > -t-1$.  In this case, we get the opposite of what we got before.  As $t$ increases, $f'=f+t$
		increases, and so we expect solutions whose slope increases as $t$ increases.  Further, after 
		at most $1$ unit of increase in $t$, we expect $f'>0$.
	\end{quote}

	In the example $f' =f+t$, there was a solution for every initial condition.  Recall that a
	solution is a \emph{function} that satisfies the differential equation and the initial conditions.
	Sometimes we can see patterns in the slope field that are not functions.

	For each of the following, draw the slope field and estimate what regions of initial conditions
	correspond to functions, and how many different ``looks'' can solutions have.
	\begin{enumerate}
		\item $y'=x/y$
		\item $y'=y/x$
		\item $y'=e^x$
		\item $y'=y(2-y)$
		\item $y'=3x^2-4$
		\item $y'=x\sqrt{y}$
		\item $y'=y^2-x$
		\item $y'=x^2-y^2$
	\end{enumerate}

	\section*{Euler's Method}

	Euler's methods is the name of a numerical algorithm to approximate the solution to
	initial value problems.  The logic follows from the idea of a linear approximation.
	If we have $y=f(x)$ and we know $y'$, we can approximate $f(x+\Delta x) \approx y+\Delta x y'=y_1$.
	Now, if we have an equation that relates $y'$ to $x$ and $y$, we can use $y_1$ to approximate
	$y_1'$, the slope of $f$ at $x+\Delta x$.  Once we have that data, we repeat.  If our initial 
	step size was small enough, we would expect to get a reasonable numerical approximation.

	\begin{definition}[Euler's Method]
		Let $y'(x)=f(x,y)$ be a first order ordinary differential equation.  The \emph{Euler approximation}
		to the initial value problem $y'=f(x,y)$ and $(x, y(x))=(x_0,y_0)$ with step size $\Delta x$ is
		the sequence of points $(x_i,y_i)$ given by $(x_0,y_0)$ if $i=0$ and 
		\[
			x_i = x_{i-1}+\Delta x
		\]
		\[
			y_i = y_{i-1} + f(x_{i-1}, y_{i-1}).
		\]
		The inductive algorithm to generate $(x_i,y_i)$ is called \emph{Euler's Method}.
	\end{definition}

	Euler's method is rather straightforward, but how do we know it actually approximates 
	solutions to the differential equation?

	\begin{definition}[Convergent Approximation]
		Let $\{(x_i,y_i)_{\Delta x}\}$ be the output of some algorithm $\mathcal A(\Delta x)$ for the initial
		value problem $y'=f(x,y)$ with initial conditions $(x,y)=(x_0,y_0)$.  We say $\mathcal A(\Delta x)$
		\emph{converges} to the solution to the initial value problem if
		\[
			\max_{(x_i,y_i)_{\Delta x}} | y_i - y(x_i)| \to 0\qquad \text{as}\qquad \Delta x\to 0
		\]
		where $y$ the true solution to the initial value problem.
	\end{definition}

	Euler's method is an algorithm with a tunable parameter, $\Delta x$, the step size.
	Saying that Euler's method converges just amounts to saying that the error
	in the estimated solution goes to zero as the step size goes to zero.

	It'd be swell if Euler's method converged.  And, in fact, it does, but not 
	necessarily all the time.

	\begin{theorem}
		Consider the initial value problem $y'=f(x,y)$ with initial conditions
		$(x_0,y_0)$.  Further suppose $\|\nabla f(x,y)\| \leq M$ for $(x,y)\in [x_0,x_0+T]\times\R$.
		Then Euler's method converges on the interval $[x_0,x_0+T]$.
	\end{theorem}

	\begin{proof}
		Let $y^*$ be the true solution to the initial value problem.  We must show that
		the Euler approximates $(x_i,y_i)$ have bounded total distance from $(x_i, y^*(x_i))$.
		
		To do this, we will consider two types of error.  Let $d_i$ be the \emph{local error}
		for $(x_i,y_i)$ and let $e_i$ be the \emph{total error} at $(x_i,y_i)$.  That is,
		let $y^*_i$ be the solution to the initial value problem $y'=f(x,y)$ with initial conditions
		$(x_i,y_i)$.  Then,
		\[
			d_i = \underbrace{y_{i-1}^*(x_{i-1}+\Delta x)}_{\text{how much
			we should have changed by}}-\qquad \underbrace{(y_{i-1} + f(x_{i-1},y_{i-1})\Delta x)}_{\text{next step}},
		\]
		and 
		\[
			e_i = \underbrace{y_i}_{\text{approximation}} -\quad \underbrace{y^*(x_i)}_{\text{exact value}}.
		\]

		If we can show $|e_i|\to 0$, we'd be done!  We will do this by bounding the local errors, $d_i$,
		and then showing that if they are small enough, $e_i$ will also be small.

		To bound $d_i$, let's use our knowledge of linear approximations.  Since $y_{i-1}^*$ is
		a solution to an initial value problem, it is differentiable.  Further, since $\nabla f$
		exists, $y_{i-1}^*$ is actually twice differentiable.  In any case,
		\[
			y_{i-1}^*(x_{i-1}+\Delta x) = y_{i-1} + {y_{i-1}^*}'\Delta x + \eta_i(\Delta x)
		\]
		\[
			=y_{i-1} + f(x_{i-1},y_{i-1})\Delta x + \eta_i(\Delta x)
		\]
		where $\eta_i(\Delta x)/\Delta x\to 0$ as $\Delta x\to 0$. Thus,
		\[
			d_i = \eta_i(\Delta x).
		\]

		Now, we'll bound the total error at the $i$th step.
		\[
			e_i = y^*(x_{i}) - y_{i} 
		\]\[=
			\underbrace{y^*(x_{i-1}) + f(x_{i-1}, y^*(x_{i-1}))\Delta x + 
			\eta^*_{i-1}(\Delta x)}_{\text{exact expantion of $y^*$}} \quad-\quad \underbrace{(y_{i-1} + f(x_{i-1}, y_{i-1})\Delta x)}_{y_i},
		\]
		where $\eta_{i-1}^*(\Delta x)$ is again an error term that goes to zero as $\Delta x\to 0$.

		Collecting terms, we now see
		\[
			e_i = y^*(x_{i-1})-y_{i-1} + \Big(f(x_{i-1}, y^*(x_{i-1})) - f(x_{i-1}, y_{i-1}) \Big)\Delta x + \eta^*_{i-1}(\Delta x)
		\]\[
			=e_{i-1} + \Big(f(x_{i-1}, y^*(x_{i-1})) - f(x_{i-1}, y_{i-1}) \Big)\Delta x + \eta^*_{i-1}(\Delta x).
		\]
		However, we supposed $\|\nabla f\| \leq M$, so 
		\[
			\big|f(x_{i-1}, y^*(x_{i-1})) - f(x_{i-1}, y_{i-1})\big| \leq M|y^*(x_{i-1}) -  y_{i-1}|
			= M|e_{i-1}|.
		\]
		Putting this all together, we see
		\[
			|e_i| \leq |e_{i-1}| + M|e_{i-1}|\Delta x + |\eta^*_{i-1}(\Delta x)|
			= (1+M\Delta x)|e_{i-1}| + |\eta^*_{i-1}(\Delta x)|.
		\]
		Let $d = \max\{|\eta_1|,|\eta_2|,\ldots,|\eta_1^*|,|\eta_2^*|,\ldots\}$.  We now have the recursive formula
		\[
			|e_i| \leq  (1+M\Delta x)|e_{i-1}| + d
		\]
		\[
			\leq (1+M\Delta x)((1+M\Delta x)|e_{i-2}| + d) + d = (1+M\Delta x)^2|e_{i-2}| + (1+(1+M\Delta x))d
		\]
		\[
			\leq (1+M\Delta x)^i|e_{0}| + \sum_{k=0}^{i-1} (1+M\Delta x)^k d.
		\]
		Since $e_0=0$ by construction, we know
		\[
			|e_i| \leq \sum_{k=0}^{i-1} (1+M\Delta x)^k d  = \frac{(1+M\Delta x)^i - 1}{M\Delta x} d.
		\]
		Now since $1+t \leq e^t$, we may replace $1+M\Delta x$ by $e^{M\Delta x}$ to finally get
		\[
			|e_i| \leq \frac{e^{iM\Delta x}-1}{M}\cdot \frac{d}{\Delta x}.
		\]
		Sine we are approximating in the interval $[x_0,x_0+T]$, we know $i\Delta x\leq T$, and so
		\[
			|e_i| \leq \frac{e^{MT}-1}{M}\cdot \frac{d}{\Delta x},
		\]
		and so $e_i$ is bounded by a constant times $d/\Delta x$.  Since $d$ was the maximum of errors
		from first order approximations whose derivatives were uniformly bounded, we know $d/\Delta x\to 0$
		as $\Delta x\to 0$.  However, we assumed something even stronger.  We assumed that
		$\nabla f$ existed and so ${y^*}''$ and ${y^*_i}''$ exist.  Thus, we actually have $d/(\Delta x^2)\to k$,
		which means the error in our approximation decreases \emph{linearly} with the size of $\Delta x$.
		


	\end{proof}

\end{document}
